\chapter{Réalisation et mise en œuvre}
\label{ch:realisation}

\section*{Introduction du chapitre}
Ce chapitre présente l'environnement de travail retenu, puis se
concentre --- conformément à la démarche itérative par fonctionnalité
décrite au chapitre~\ref{ch:contexte} --- sur les mécanismes qui
\emph{dépassent} le CRUD simple et qui constituent la valeur ajoutée
technique réelle du projet.

\section{Environnement de travail}

\begin{table}[H]
\centering
\begin{tabular}{@{}>{\bfseries}p{3.5cm}p{10.5cm}@{}}
\toprule
Composant & Version / outil \\
\midrule
Backend & Spring Boot 3.4.5, Java 17+, Maven \\
Frontend & Next.js 15 (App Router, Turbopack), TypeScript strict, Tailwind v4 \\
Base de données & PostgreSQL 16, migrations versionnées Flyway \\
Cache / verrous & Redis 7 \\
Conteneurisation & Docker, Docker Compose \\
Versionnement & Git / GitHub \\
Modélisation UML & PlantUML (diagrammes de classes, cas d'utilisation,
  séquence), Python/\texttt{mingrammer} (architecture, déploiement) \\
Tests & JUnit 5 + Mockito (backend), Vitest + React Testing Library
  (frontend), Playwright (bout en bout) \\
\bottomrule
\end{tabular}
\caption{Environnement logiciel et outils du projet}
\label{tab:techno-outils}
\end{table}
\tabdesc{Chaque outil retenu répond à une contrainte du projet plutôt qu'à
une préférence personnelle ; le Tableau~\ref{tab:stack} du chapitre
précédent détaille la justification des choix technologiques
d'architecture proprement dits.}

\section{Vue d'ensemble des modules implémentés}

Le backend expose 29 contrôleurs REST couvrant l'ensemble du cycle de vie
pédagogique : authentification (avec 2FA optionnelle), catalogue de cours
et éditeur de contenu par blocs, système de quiz à quatre portées
(section, module, UF, année), gestion de stage avec signature externe,
certification, messagerie, notifications, campagnes email, bourse à
l'emploi, recherche globale et outils d'administration (utilisateurs,
groupes, journal d'audit, paramètres système). Plutôt que de décrire
exhaustivement chaque endpoint CRUD --- déjà couvert au
chapitre~\ref{ch:analyse-besoins} --- ce chapitre se concentre sur les
mécanismes qui \emph{dépassent} le CRUD simple et qui constituent la
valeur ajoutée technique réelle du projet.

\section{Points forts techniques et valeur ajoutée}
\label{sec:points-forts}

Les onze points suivants sont classés par ordre d'impact décroissant :
du plus démonstratif d'une compétence technique ou d'une valeur ajoutée
utilisateur, au plus secondaire. Chaque affirmation renvoie à un fichier
et, quand c'est utile, à une ligne précise du code source.

%-----------------------------------------------------------------------
\subsection{Génération de questions de quiz par intelligence artificielle, avec repli entre fournisseurs}
\label{sec:generation-ia}

\textbf{Contexte du problème.} Rédiger manuellement des questions de quiz
pour chaque leçon est chronophage pour l'équipe pédagogique. Un service
d'IA générative peut accélérer ce travail, mais dépendre d'un seul
fournisseur externe (quota, panne, latence) rendrait la fonctionnalité peu
fiable en production.

\textbf{Solution apportée.} \texttt{QuestionGenerationService}
(\texttt{service/QuestionGenerationService.java}) interroge d'abord l'API
Gemini (\texttt{callGemini}, ligne 120) ; en cas d'échec --- réseau, réponse
HTTP non-2xx, JSON invalide, ou aucune question extraite --- il bascule
automatiquement vers l'API Grok (\texttt{callGrok}, ligne 140), avant
d'échouer explicitement si les deux fournisseurs sont indisponibles. Les
deux méthodes sont volontairement \texttt{package-private} (et non
\texttt{private}) pour être simulées (\texttt{Mockito.spy()}) dans
\texttt{QuestionGenerationServiceTest} sans appel réseau réel. Le résultat
brut est nettoyé (extraction hors des blocs \texttt{```json}), parsé, puis
chaque question générée passe par \emph{exactement} la même validation
(\texttt{addQuestion}) qu'une question créée manuellement --- une question
invalide dans le lot est rapportée comme erreur individuelle sans annuler
les autres. L'appel est en outre protégé par
\texttt{RateLimitService.checkAiGenerationAllowed} (20 requêtes / 5 minutes
par utilisateur) pour contenir le coût d'un appel à une API payante.

\begin{figure}[H]
  \centering
  \includegraphics[width=\textwidth,height=0.85\textheight,keepaspectratio]{diagramme_sequence_generation_questions_ia.png}
  \caption{Génération de questions par IA avec repli Gemini vers Grok}
  \label{fig:seq-generation-ia}
\end{figure}

Côté interface, cette mécanique reste volontairement discrète : le
formateur travaille depuis la banque de quiz habituelle
(Figure~\ref{fig:capture-banque-quiz}) et déclenche la génération par une
seule icône, sans quitter l'écran d'édition des questions.

\figimg{screenshots/capture_banque_quiz.png}{0.95}
  {Interface formateur --- banque de quiz}{capture-banque-quiz}
  {Chaque question est affichée avec ses options et la bonne réponse
   surlignée en vert ; l'icône en forme d'étincelle (en haut à droite du
   quiz sélectionné) ouvre la modale de génération par IA
   (Figure~\ref{fig:capture-generation-ia}).}

\figimg{screenshots/capture_generation_ia.png}{0.75}
  {Modale de génération de questions par IA}{capture-generation-ia}
  {Le formateur colle un extrait de cours, choisit le type de question et
   le nombre à générer ; les questions produites arrivent directement
   dans la banque, prêtes à être relues avant publication.}

\textbf{Bénéfice concret.} Une simple intégration à un seul fournisseur
aurait suffi pour une démonstration, mais pas pour un usage réel en
production, où une panne ou un changement de quota chez Gemini aurait
purement et simplement désactivé la fonctionnalité. Le mécanisme de repli
et la réutilisation de la validation existante garantissent que le contenu
généré par IA respecte les mêmes invariants qu'un contenu créé à la main,
sans code de persistance dupliqué.

%-----------------------------------------------------------------------
\subsection{Correction hybride des quiz : automatique et manuelle, avec réconciliation d'état}
\label{sec:correction-hybride}

\textbf{Contexte du problème.} Un même quiz peut mélanger des questions à
choix (auto-corrigeables) et des questions à réponse libre (nécessitant un
jugement humain). Une solution naïve traiterait chaque soumission comme
« terminée » dès l'envoi, ce qui est incorrect si une partie de la note
dépend encore d'une correction à venir.

\textbf{Solution apportée.} \texttt{QuizAttemptService.submit()}
(\texttt{service/QuizAttemptService.java}) calcule le score des questions
auto-corrigeables par égalité stricte d'ensembles d'identifiants
d'options sélectionnées --- ce qui signifie qu'un choix multiple ne donne
\emph{aucun} crédit partiel : une réponse presque correcte compte comme
fausse. Si le quiz contient au moins une question \texttt{ESSAY}, la
tentative passe au statut \texttt{PENDING\_REVIEW} plutôt que
\texttt{PASSED}/\texttt{FAILED}. \texttt{QuizGradingService.gradeEssay()}
persiste chaque note manuelle, puis \texttt{finalizeIfFullyGraded()}
vérifie si \emph{toutes} les questions à réponse libre de la tentative ont
désormais une note --- ce n'est qu'à ce moment que le score final est
recalculé et que la tentative bascule vers son statut définitif, en
réutilisant \texttt{applyPassEffects()} --- la même méthode qui gère
badges, notifications et éligibilité au certificat pour une soumission
entièrement automatique. Une autre difficulté résolue : le démarrage
concurrent d'une tentative est sérialisé par un verrou Postgres
(\texttt{pg\_advisory\_xact\_lock}) pour empêcher qu'un double clic ne crée
deux tentatives \texttt{IN\_PROGRESS} simultanées --- comportement vérifié
par un script de test de charge (\texttt{loadtest/quiz-start-concurrency.js}).

\begin{figure}[H]
  \centering
  \includegraphics[width=\textwidth,height=0.85\textheight,keepaspectratio]{diagramme_sequence_correction_quiz_hybride.png}
  \caption{Correction hybride d'un quiz et réconciliation d'état}
  \label{fig:seq-correction-hybride}
\end{figure}

\figimg{screenshots/capture_quiz_passation.png}{0.7}
  {Passation d'un quiz --- côté apprenant}{capture-quiz-passation}
  {Chronomètre, navigation question par question et libellé du quiz sont
   affichés en continu ; la soumission finale déclenche côté serveur la
   logique de correction hybride décrite ci-dessus.}

\textbf{Bénéfice concret.} Sans cette réconciliation explicite, un
formateur pourrait involontairement déclencher une notification
« quiz réussi » ou l'émission d'un certificat avant qu'une réponse libre
ait été relue --- une erreur silencieuse difficile à détecter en
production. La réutilisation de \texttt{applyPassEffects()} garantit aussi
qu'un apprenant reçoit exactement le même traitement (badge, certificat)
qu'il ait été corrigé automatiquement ou manuellement.

%-----------------------------------------------------------------------
\subsection{Progression pédagogique séquentielle par UF, avec cascade vers l'éligibilité au certificat}
\label{sec:progression-uf}

\textbf{Contexte du problème.} Le cursus impose un ordre : un apprenant ne
doit pas pouvoir accéder au contenu d'une UF avant d'avoir complété la
précédente, ni obtenir un certificat sans avoir réellement terminé
l'intégralité du parcours --- y compris les étapes qui exigent un jugement
humain (validation du stage par le Directeur).

\textbf{Solution apportée.} \texttt{ProgressionService.isUfFullyDone()}
(voir aussi la note du diagramme de classes,
Figure~\ref{fig:classes-pedagogie}) combine trois conditions : le contenu de tous les
modules de l'UF est complété, le quiz de fin d'UF est réussi \emph{si} il
existe (sinon la condition est neutre), et
\texttt{UfValidationService.isValidated()} confirme une validation manuelle
--- mais uniquement pour les deux UF qui l'exigent réellement (\texttt{UF 5}
= stage, \texttt{UF 11} = soutenance ; toutes les autres UF renvoient
\texttt{true} sans intervention humaine). Les apprenants rattachés à un
groupe (cohorte présentielle/hybride) échappent à cette règle séquentielle
automatique : leur accès est piloté explicitement par
\texttt{GroupContentAssignment}, avec des dates d'ouverture fixées par le
Directeur --- deux stratégies de déblocage cohabitant dans le même modèle
selon le profil d'apprenant. En bout de chaîne,
\texttt{CertificateService.tryIssueIfEligible()} ne délivre un certificat
que si \emph{tous} les modules de la formation sont complétés --- la
cascade UF~$\rightarrow$~module~$\rightarrow$~certificat n'est jamais
court-circuitée.

\begin{figure}[H]
  \centering
  \includegraphics[width=\textwidth,height=0.85\textheight,keepaspectratio]{diagramme_sequence_progression_uf_certificat.png}
  \caption{Progression séquentielle par UF et cascade vers le certificat}
  \label{fig:seq-progression-uf}
\end{figure}

\textbf{Bénéfice concret.} Un modèle plus simple (un simple pourcentage de
progression global) n'aurait pas pu représenter la contrainte réelle du
Directeur : certaines étapes ne se valident pas automatiquement, quel que
soit le score obtenu. La double stratégie individuel/groupe permet en
outre au même système de servir aussi bien les apprenants en ligne
autonomes que les cohortes encadrées en présentiel.

%-----------------------------------------------------------------------
\subsection{Contrôle d'accès à granularité fine, au-delà du RBAC}
\label{sec:acces-proprietaire}

\textbf{Contexte du problème.} Un rôle (\texttt{@PreAuthorize}) suffit à
dire \emph{qu'un} apprenant peut lire ses messages, pas \emph{lequel}. Sans
vérification supplémentaire, un utilisateur authentifié pourrait accéder
aux conversations, notifications ou notes personnelles d'un autre
utilisateur simplement en devinant un identifiant dans l'URL.

\textbf{Solution apportée.} Plusieurs services implémentent un contrôle de
propriété explicite, systématiquement revérifié à \emph{chaque} opération
(pas seulement à l'affichage de la liste) : \texttt{MessagingService.assertAccess()}
distingue les conversations directes (vérification d'une ligne
\texttt{ConversationParticipant}) des salons de cohorte (vérification du
groupe partagé, avec accès de plein droit pour le staff) ;
\texttt{NotificationService.markRead()} et
\texttt{LessonNoteService.requireOwnNote()} comparent directement le
propriétaire de la ressource à l'utilisateur authentifié ;
\texttt{StageService} va plus loin avec trois méthodes d'assertion
distinctes selon le type de document (document réservé au Directeur,
document réservable à l'apprenant concerné uniquement, avec en plus une
fenêtre de dates pour l'upload).

\begin{figure}[H]
  \centering
  \includegraphics[width=\textwidth,height=0.85\textheight,keepaspectratio]{diagramme_sequence_controle_acces_proprietaire.png}
  \caption{Contrôle d'accès à granularité fine (exemple de la messagerie)}
  \label{fig:seq-acces-proprietaire}
\end{figure}

\textbf{Bénéfice concret.} Ce pattern --- vérification de propriété au
niveau service, jamais seulement au niveau contrôleur --- ferme une classe
entière de vulnérabilités de type IDOR (\emph{Insecure Direct Object
Reference}, OWASP Top 10~\cite{owasp-top10}). Sa mise en \oe uvre n'est
toutefois pas uniforme dans tout le code : voir la limite identifiée au
\S\ref{sec:limites}.

%-----------------------------------------------------------------------
\subsection{Limitation de débit Redis, appliquée de façon transverse}
\label{sec:rate-limiting}

\textbf{Contexte du problème.} Plusieurs endpoints exposent une surface
d'abus possible : tentatives de connexion (force brute), appels à des API
IA payantes, envoi de messages (spam), upload de fichiers.

\textbf{Solution apportée.} \texttt{RateLimitService} implémente un
compteur à fenêtre fixe avec exactement deux primitives Redis --- \texttt{INCR}
puis un \texttt{EXPIRE} conditionnel posé uniquement à la création du
compteur --- exposé via huit méthodes dédiées
(\texttt{checkLoginAllowed}, \texttt{checkAiGenerationAllowed},
\texttt{checkMessageAllowed}, etc.), chacune avec son propre seuil et sa
propre fenêtre temporelle. Le dépassement lève une
\texttt{RateLimitException} annotée \texttt{@ResponseStatus(TOO\_MANY\_REQUESTS)},
traduite automatiquement en HTTP~429 sans gestionnaire d'exception dédié.

\textbf{Bénéfice concret.} Un seul mécanisme réutilisable protège à la
fois la sécurité (force brute sur la connexion) et le budget
d'exploitation (appels IA facturés à l'usage) --- sans dépendance
supplémentaire, Redis étant déjà présent dans la stack pour d'autres
besoins.

%-----------------------------------------------------------------------
\subsection{Certificats vérifiables publiquement et partage LinkedIn}
\label{sec:certificat-linkedin}

\textbf{Contexte du problème.} Un certificat PDF classique n'offre aucun
moyen à un tiers (recruteur, partenaire) de vérifier son authenticité, et
ne valorise pas la réussite de l'apprenant sur les réseaux professionnels.

\textbf{Solution apportée.} \texttt{CertificateService} génère un code de
vérification aléatoire (UUID tronqué, 16 caractères) exposé par un
endpoint public \emph{sans authentification}
(\texttt{GET /api/certificates/verify/\{code\}}), qui ne renvoie que des
informations déjà publiques par nature (nom de l'apprenant, formation,
date --- jamais d'email ni de PDF téléchargeable). Côté frontend, la page
Next.js \texttt{/verify/[code]} est un \emph{composant serveur} dont la
fonction \texttt{generateMetadata} interroge cette API pour produire les
balises Open Graph directement dans le HTML initial --- condition
nécessaire pour que le robot d'indexation de LinkedIn, qui n'exécute pas
de JavaScript, affiche un aperçu correct lors du partage.

\figimg{screenshots/capture_verification_certificat.png}{0.7}
  {Page publique de vérification d'un certificat}{capture-verification}
  {Accessible sans authentification à l'adresse \texttt{/verify/[code]},
   cette page ne révèle que des informations déjà publiques (nom,
   formation, date) et propose un bouton de partage direct vers
   LinkedIn.}

\textbf{Bénéfice concret.} Cette fonctionnalité transforme un document
statique en preuve vérifiable et en outil de visibilité professionnelle
pour l'apprenant --- une valeur ajoutée directement alignée sur la mission
de l'académie (employabilité).

%-----------------------------------------------------------------------
\subsection{Recherche globale multi-source, avec confidentialité des questions de quiz}
\label{sec:recherche-globale}

\textbf{Contexte du problème.} Une recherche naïve sur les seuls titres de
modules/leçons est peu utile. L'étendre au contenu réel (texte, PDF,
questions de quiz) introduit un risque : un apprenant pourrait utiliser la
recherche pour retrouver directement l'énoncé --- voire la réponse --- d'une
question d'examen.

\textbf{Solution apportée.} \texttt{SearchService} interroge cinq sources
(titres de modules, titres de leçons, corps de texte des blocs \texttt{TEXT},
texte extrait des blocs \texttt{PDF}, prompts de questions), avec un
extrait (\emph{snippet}) calculé autour du premier terme trouvé.
\texttt{SearchController} détermine \texttt{includeQuestions} à partir du
rôle du principal authentifié (\texttt{isStaff()}) \emph{avant} d'appeler le
service : si l'appelant n'est pas membre du staff, le dépôt de questions
n'est \textbf{jamais interrogé} --- la confidentialité n'est pas un simple
filtrage d'affichage après coup. L'extraction du texte des PDF a lieu au
moment de l'upload (\texttt{MediaService}, Apache PDFBox), en échec
silencieux (best-effort) pour ne jamais bloquer un import de fichier.

\textbf{Bénéfice concret.} La recherche devient un outil pédagogique
(retrouver un passage de cours) sans devenir un outil de triche --- un
compromis qui aurait été manqué par une implémentation qui se contenterait
de masquer les résultats côté interface.

%-----------------------------------------------------------------------
\subsection{Bourse à l'emploi dérivée de la possession d'un certificat}
\label{sec:bourse-emploi}

\textbf{Contexte du problème.} L'accès à la bourse à l'emploi doit être
réservé aux diplômés --- une notion qui pourrait être modélisée par un
simple champ booléen \texttt{isAlumni} sur \texttt{User}, au risque de se
désynchroniser de la réalité (diplôme révoqué, erreur manuelle).

\textbf{Solution apportée.} \texttt{JobOfferService.listForLearner()} ne
stocke aucun champ dédié : le statut de diplômé est \emph{dérivé} de
l'existence d'une ligne \texttt{Certificate} pour cet utilisateur et cette
formation (\texttt{certificateRepository.findByUserIdAndFormationId}),
exactement la même requête que celle utilisée pour afficher le certificat
lui-même. Un apprenant non diplômé reçoit une erreur 403 explicite
(« Réservé aux diplômés »). Le staff, lui, passe par une route
\texttt{/admin/job-offers} entièrement distincte qui ne traverse jamais
cette vérification.

\figimg{screenshots/capture_bourse_emploi.png}{0.95}
  {Espace alumni --- bourse à l'emploi}{capture-bourse-emploi}
  {Vue telle qu'affichée à un apprenant diplômé : offres réelles
   d'entreprises partenaires du secteur, avec lien de candidature et
   contact direct par courriel pour chacune.}

\textbf{Bénéfice concret.} L'état « diplômé » ne peut jamais être
incohérent avec la réalité du certificat --- il n'existe qu'un seul endroit
où cette vérité est stockée (principe \emph{single source of truth}).

%-----------------------------------------------------------------------
\subsection{Envoi de campagnes email asynchrone à débit maîtrisé}
\label{sec:campagne-email}

\textbf{Contexte du problème.} Envoyer plusieurs centaines d'emails de
façon synchrone bloquerait la requête HTTP de l'administrateur, et un
envoi trop rapide risquerait de faire percevoir le serveur SMTP comme
source de spam.

\textbf{Solution apportée.} \texttt{EmailCampaignService.send()} persiste
immédiatement tous les destinataires en statut \texttt{PENDING}, bascule la
campagne à \texttt{SENDING}, puis délègue à
\texttt{CampaignSenderWorker.sendAsync()} --- une classe séparée
spécifiquement parce qu'une méthode \texttt{@Async} appelée depuis la même
classe contournerait le proxy AOP de Spring et ne serait jamais réellement
asynchrone. L'envoi traite les destinataires un par un avec une pause fixe
de 200~ms, sur un pool de threads volontairement restreint (2 à 4 threads)
pour ne jamais saturer le serveur SMTP. Chaque échec individuel est
capturé, journalisé et marqué \texttt{FAILED} sans interrompre l'envoi des
destinataires suivants.

\textbf{Bénéfice concret.} L'administrateur récupère l'interface
immédiatement après avoir lancé l'envoi, et l'état de chaque destinataire
(en base, pas dans une file séparée) reste consultable et survit à un
redémarrage du serveur --- avec une limite assumée décrite au
\S\ref{sec:limites}.

%-----------------------------------------------------------------------
\subsection{Détection d'inactivité (alerte précoce)}
\label{sec:early-warning}

\textbf{Contexte du problème.} Repérer les apprenants décrocheurs
permet une intervention pédagogique avant l'échec.

\textbf{Solution apportée.} \texttt{EarlyWarningService.findInactiveStudents(days)}
calcule, pour chaque apprenant, le nombre de jours écoulés depuis la
dernière mise à jour de sa \texttt{LessonProgress} --- avec repli sur la
date de création du compte si aucune progression n'existe encore --- et
retourne la liste triée par durée d'inactivité décroissante, consultée à
la demande depuis le tableau de bord analytique (aucune tâche planifiée,
aucun effet de bord automatique).

\textbf{Bénéfice concret.} Un signal simple mais fiable, calculé à partir
d'une donnée déjà existante plutôt que d'un nouveau mécanisme de suivi
dédié.

%-----------------------------------------------------------------------
\subsection{Verrou d'édition collaboratif (distinct de la progression séquentielle)}
\label{sec:verrou-edition}

\textbf{Contexte du problème.} Deux membres du staff éditant la même leçon
en même temps risqueraient d'écraser mutuellement leurs modifications.

\textbf{Solution apportée.} \texttt{LessonLockService} pose un verrou
Redis nominatif (\texttt{lesson\_lock:\{lessonId\}}, TTL 30 minutes,
renouvelable) identifiant qui édite actuellement une leçon ; toute
tentative concurrente reçoit une \texttt{LessonLockedException} (HTTP 409)
nommant la personne qui détient le verrou. \emph{Précision utile pour ce
rapport} : ce mécanisme ne doit pas être confondu avec le verrouillage
séquentiel du parcours apprenant (\S\ref{sec:progression-uf}, implémenté
dans \texttt{ProgressionService}) --- il s'agit d'une garde d'édition
collaborative côté back-office, sans rapport avec la progression
pédagogique.

\textbf{Bénéfice concret.} Un conflit d'édition est prévenu plutôt que
constaté après coup, avec un message explicite plutôt qu'un écrasement
silencieux.

\subsection{Tableau de synthèse}

Le Tableau~\ref{tab:synthese-points-forts} regroupe, dans l'ordre d'impact
retenu tout au long de cette section, les onze mécanismes présentés :
leur emplacement précis dans le code source, et l'argument de valeur
ajoutée qui justifie leur mise en avant --- réutilisable tel quel comme
support de soutenance.

\begin{table}[H]
  \centering
  \small
  \caption{Synthèse des points forts techniques et de leur valeur ajoutée}
  \label{tab:synthese-points-forts}
  \begin{tabular}{@{}p{2.6cm}p{4.1cm}p{6.3cm}@{}}
    \toprule
    Mécanisme & Localisation & Valeur ajoutée \\
    \midrule
    Génération de questions par IA (repli Gemini/Grok) & \texttt{Question\allowbreak Generation\allowbreak Service} & Fonctionnalité IA fiable en production, jamais dépendante d'un seul fournisseur externe. \\
    Correction hybride de quiz (auto + manuelle) & \texttt{QuizAttempt\allowbreak Service}, \texttt{QuizGrading\allowbreak Service} & Un score final toujours cohérent, jamais figé avant qu'une correction manuelle en attente ne soit terminée. \\
    Progression séquentielle par UF $\rightarrow$ certificat & \texttt{Progression\allowbreak Service}, \texttt{UfValidation\allowbreak Service}, \texttt{Certificate\allowbreak Service} & Respecte les validations humaines obligatoires (stage, soutenance) sans jamais les court-circuiter automatiquement. \\
    Contrôle d'accès à granularité fine & \texttt{Messaging\allowbreak Service}, \texttt{Notification\allowbreak Service}, \texttt{LessonNote\allowbreak Service}, \texttt{Stage\allowbreak Service} & Ferme une classe de vulnérabilités IDOR que le seul RBAC ne peut pas couvrir. \\
    Limitation de débit Redis transverse & \texttt{RateLimit\allowbreak Service} & Protège à la fois la sécurité (force brute) et le budget d'exploitation (appels IA payants) avec un seul mécanisme réutilisable. \\
    Certificat vérifiable publiquement + partage LinkedIn & \texttt{Certificate\allowbreak Service}, page \texttt{/verify/[code]} & Transforme un diplôme statique en preuve vérifiable et en outil d'employabilité réelle. \\
    Recherche globale avec confidentialité des questions & \texttt{Search\allowbreak Service}, \texttt{Search\allowbreak Controller} & Un outil pédagogique qui ne peut jamais devenir un outil de triche. \\
    Bourse à l'emploi dérivée du certificat & \texttt{JobOffer\allowbreak Service} & Le statut « diplômé » ne peut jamais se désynchroniser de la réalité (source de vérité unique). \\
    Campagnes email asynchrones à débit maîtrisé & \texttt{EmailCampaign\allowbreak Service}, \texttt{CampaignSender\allowbreak Worker} & L'administrateur n'attend jamais l'envoi complet, sans jamais saturer le serveur SMTP. \\
    Détection d'inactivité & \texttt{EarlyWarning\allowbreak Service} & Signal fiable pour l'intervention pédagogique, sans mécanisme de suivi dédié supplémentaire. \\
    Verrou d'édition collaboratif & \texttt{LessonLock\allowbreak Service} & Prévient un conflit d'édition entre formateurs plutôt que de le constater après coup. \\
    \bottomrule
  \end{tabular}
\end{table}

\section{Limites identifiées et pistes d'amélioration}
\label{sec:limites}

Deux limites concrètes, identifiées en analysant le code plutôt que
supposées, sont assumées ici plutôt que passées sous silence :

\begin{itemize}
  \item \textbf{Incohérence du contrôle de propriété.} Contrairement aux
    services décrits au \S\ref{sec:acces-proprietaire},
    \texttt{AssignmentService.gradeSubmission()} et
    \texttt{AssignmentService.delete()} ne vérifient aucune propriété de
    ressource au-delà du rôle \texttt{@PreAuthorize("hasAnyRole('ADMIN','FORMATEUR')")} :
    un formateur peut ainsi noter ou supprimer un devoir rattaché à un
    module qu'il n'encadre pas. \emph{Piste d'amélioration} : appliquer le
    même pattern \texttt{assertAccess} que dans \texttt{MessagingService}.
  \item \textbf{Reprise après incident non automatisée.} L'état d'une
    campagne email interrompue par un redémarrage serveur (destinataires
    encore \texttt{PENDING}) survit en base de données, mais aucun
    mécanisme planifié ne relance automatiquement l'envoi --- une reprise
    manuelle est nécessaire. \emph{Piste d'amélioration} : un
    \texttt{@Scheduled} au démarrage de l'application qui détecte les
    campagnes \texttt{SENDING} avec des destinataires \texttt{PENDING} et
    relance \texttt{sendAsync()}.
\end{itemize}

\section*{Conclusion du chapitre}
Ce chapitre a détaillé, par ordre d'impact, les mécanismes qui distinguent
\Entreprise{} d'une simple application CRUD : intégration IA résiliente,
correction hybride avec réconciliation d'état, progression métier en
cascade, contrôle d'accès à granularité fine, et les limites identifiées
en toute transparence. Le chapitre~\ref{ch:tests-deploiement-bilan} présente
maintenant la démarche de validation, de déploiement et le bilan du stage.
